%!TEX root = ../main.tex
\section{Prescriptive Data Analytics of COVID-19}
\label{sec:prescriptive}


%\stitle{Ad-hoc Visualizations}.
We also design ad-hoc visualizations to answer specific questions. In terms of COVID-19, besides publicly available datasets, we also have private trajectory data of potentially infected persons. Based on which we have designed two map-based visualizations, one to show infection paths of these patients (see Figure~\ref{fig:infection}), and the other to show the level of risk for each area (see Figure~\ref{fig:heatmap}) and thus suggest the authorities to take different anti-epidemic policies for different areas.

%\add{You should merge ad-hoc visualizations to Section 6.}

%%%%%%%%%%%%%%%%%%%%%%%%%%%%%%%%%%%%
\begin{figure}[t!]
	\vspace{-.5em}
	\centering
	\includegraphics[width=.65\columnwidth]{figs/InfectionPath.pdf} 
	\caption{Tracking Infection Path}
	\label{fig:infection}
\end{figure}
%%%%%%%%%%%%%%%%%%%%%%%%%%%%%%%%%%%%

\subsection{Infection Path}
Based on the trajectory data of (potentially) infected persons, we can support to visualize and track the infection path at the high-level. Taking Figure~\ref{fig:infection} as an example, a person started from {\em Beijing Haidian Hospital} at 13:28 on 2020-02-28, and walked through several streets and finally arrived at his/her neighborhood. Therefore, we cluster those trajectories of infected persons to find a group of high-risk roads. 
%
Moreover, we devise a trajectory similarity search technique to find other trajectories similar to those trajectories of infected persons. It will help us to find and report the potentially high-risk groups. 
Thus, the authorities can take different anti-epidemic policies against people at different risk levels.

\stitle{Findings and Insights.}
According to the sample trajectory data of (potentially) infected persons provided by China Mobile Limited, we find that most of the trajectories passed through some living places such as restaurants and supermarkets, and finally return home. Some trajectories intersect with public transportation such as trains and subways. There are a few trajectories with several other trajectories that coincide, indicating that they may be traveling together.
%
Since many trajectories have visited supermarkets and other living places, the government can suggest that people go to supermarkets as little as possible, or that different people go to supermarkets at different times, and limits the number of customers in supermarkets.

\subsection{High-risk Areas Discovery}
To further explore the trajectory data of (potentially) infected persons, we also design a visualization for the trajectory points using a heatmap. 
Figure~\ref{fig:heatmap} gives an at-a-glance understanding of the spatial distribution of the potentially infected persons from Hubei, China to Beijing, China.
Those ``hot'' areas mean there have more potentially infected persons visit. 
We also indicate those neighborhoods that have several confirmed cases in the heatmap by the symbol~\includegraphics[height=1.\baselineskip]{figs/marker.png}.

%
%Take Figure~\ref{fig:community} for example, users can understand the COVID-19 situations near {\em Tsinghua University, Beijing, China} by a location search box. Note that this module only supports for the mainland China area currently.


\stitle{Findings and Insights.}
According to the location data of potentially infected persons from Hubei to Beijing provided by China Mobile Limited, we make the following observations:
(1) Most of the potential persons arrive in Beijing through transportation hubs such as railway stations (\includegraphics[height=0.8\baselineskip]{figs/train.eps} in Figure~\ref{fig:heatmap}) and airports. Intuitively, such transportation hubs likely play a significant role in the spread of the coronavirus in Beijing City. Therefore, the government should take corresponding measures for these transportation hubs to minimize the risk of the transmission of coronavirus;
(2) Many potentially infected persons visit some commercial places (\eg supermarkets) and end up staying in their residence; and
(3) There is a high overlap between areas where high-risk people are active and neighborhoods containing confirmed cases. 
%Therefore, we can infer that other places visited by high-risk groups may also be dangerous.
Intuitively, we can infer that other places (\eg the red rectangles) visited by high-risk groups may also be dangerous. 
Therefore, the authorities can take different anti-epidemic policies against people at different risk levels.
Note that, for those ``hot'' areas do not appear confirmed cases (\eg the red rectangles), some of them are later reported confirmed cases by the government.
%
Thus, the authorities can take precautions against these areas in advance and take different anti-epidemic policies against different areas to achieve refined and effective management.

%%%%%%%%%%%%%%%%%%%%%%%%%%%%%%%%%%%%
\begin{figure}[t!]
	\centering
	\includegraphics[width=.7\columnwidth]{figs/phone.pdf}
	\caption{High-risk Area Discovery}
	\label{fig:heatmap}
\end{figure}
%%%%%%%%%%%%%%%%%%%%%%%%%%%%%%%%%%%%





%%%%%%%%%%%%%%%%%%%%%%%%%%%%%%%%%%%%
%\begin{figure}[t!]
%	\centering
%	\includegraphics[width=.5\columnwidth]{figs/community.pdf}
%	\caption{Find Confirmed Cases Around You}
%	\label{fig:community}
%\end{figure}
%%%%%%%%%%%%%%%%%%%%%%%%%%%%%%%%%%%

%%%%%%%%%%%%%%%%%%%%%%%%%%%%%%%%%%%%%
%\begin{figure}[t!]
%	\centering
%	\includegraphics[width=.7\columnwidth]{figs/similarity_trend1.pdf}
%	\caption{Similar Trend of Confirmed Cases}
%	\label{fig:similar_trend}
%\end{figure}
%%%%%%%%%%%%%%%%%%%%%%%%%%%%%%%%%%%%\\

%\stitle{Any other story you can tell from China mobile data}